\section{AMC 12 12A 2005}
        \begin{enumerate}
        	\item (\textit{AMC 12 12A 2005}) Two is $10 \%$ of $x$ and $20 \%$ of $y$. What is $x - y$?


 $\textbf{(A) } 1 \qquad \textbf{(B) } 2 \qquad \textbf{(C) } 5 \qquad \textbf{(D) } 10 \qquad \textbf{(E) } 20$

 \textbf{Answer: }D



\textbf{Solution}

$2 = \frac {1}{10}x \Longrightarrow x = 20,\quad 2 = \frac{1}{5}y \Longrightarrow y = 10,\quad x-y = 20 - 10=10  \mathrm{(D)}$.

 


	\item (\textit{AMC 12 12A 2005}) The equations $2x + 7 = 3$ and $bx - 10 = - 2$ have the same solution for $x$. What is the value of $b$?


 $\textbf{(A) } -8 \qquad \textbf{(B) } -4 \qquad \textbf{(C) } -2 \qquad \textbf{(D) } 4 \qquad \textbf{(E) } 8$

 \textbf{Answer: }B



\textbf{Solution}

$2x + 7 = 3 \iff x = -2$, so we require $-2b-10 = -2 \iff -2b = 8 \iff b = \boxed{\textbf{(B) } -4}$.

 


	\item (\textit{AMC 12 12A 2005}) A rectangle with a diagonal of length $x$ is twice as long as it is wide. What is the area of the rectangle?


 $\textbf{(A) } \frac{1}{4}x^2 \qquad \textbf{(B) } \frac{2}{5}x^2 \qquad \textbf{(C) } \frac{1}{2}x^2 \qquad \textbf{(D) } x^2 \qquad \textbf{(E) } \frac{3}{2}x^2$

 \textbf{Answer: }B



\textbf{Solution}

Let $w$ be the width, so the length is $2w$. By the \href{https://artofproblemsolving.com/wiki/index.php?title=Pythagorean\_Theorem}{Pythagorean Theorem}, $w^2 + 4w^2 = x^2 \Longrightarrow \frac{x}{\sqrt{5}} = w$. The area of the rectangle is $(w)(2w) = 2w^2 = 2\left(\frac{x}{\sqrt{5}}\right)^2 = \frac{2}{5}x^2 \ \mathrm{(B)}$.

 


	\item (\textit{AMC 12 12A 2005}) A store normally sells windows at $\$100$ each. This week the store is offering one free window for each purchase of four. Dave needs seven windows and Doug needs eight windows. How many dollars will they save if they purchase the windows together rather than separately?


 $\textbf{(A) } 100 \qquad \textbf{(B) } 200 \qquad \textbf{(C) } 300 \qquad \textbf{(D) } 400 \qquad \textbf{(E) } 500$

 \textbf{Answer: }A



\textbf{Solution}

For $n$ windows, the store offers a discount of $100 \cdot \left\lfloor\frac{n}{5}\right\rfloor$ (\href{https://artofproblemsolving.com/wiki/index.php?title=Floor\_function}{floor function}). Dave receives a discount of $100 \cdot \left\lfloor \frac{7}{5}\right \rfloor = 100$ and Doug receives a discount of $100 \cdot \left\lfloor \frac{8}{5}\right\rfloor = 100$. These amount to $200$ dollars in discounts. Together, they receive a discount of $100 \cdot \left\lfloor \frac{15}{5} \right\rfloor = 300$, so they save $300-200=100\ \mathrm{(A)}$.

 


	\item (\textit{AMC 12 12A 2005}) The average (mean) of $20$ numbers is $30$, and the average of $30$ other numbers is $20$. What is the average of all $50$ numbers?


 $\textbf{(A) } 23 \qquad \textbf{(B) } 24 \qquad \textbf{(C) } 25 \qquad \textbf{(D) } 26 \qquad \textbf{(E) } 27$

 \textbf{Answer: }B



\textbf{Solution 1}

The sum of the first $20$ numbers is $20 \cdot 30$ and the sum of the other $30$ numbers is $30\cdot 20$. Hence the overall average is $\frac{20 \cdot 30 + 30 \cdot 20}{50} = 24 \ \mathrm{(B)}$.

 



\textbf{Solution 2}

This is just the harmonic mean, so the answer is $\frac{2 \cdot 20 \cdot 30}{20+30}=24 \ \mathrm{(B)}$.

 Solution by franzliszt

 


	\item (\textit{AMC 12 12A 2005}) Josh and Mike live $13$ miles apart. Yesterday Josh started to ride his bicycle toward Mike's house. A little later Mike started to ride his bicycle toward Josh's house. When they met, Josh had ridden for twice the length of time as Mike and at four-fifths of Mike's rate. How many miles had Mike ridden when they met?


 $\textbf{(A) } 4 \qquad \textbf{(B) } 5 \qquad \textbf{(C) } 6 \qquad \textbf{(D) } 7 \qquad \textbf{(E) } 8$

 \textbf{Answer: }B



\textbf{Solution}

Let $D_J, D_M$ be the distances traveled by Josh and Mike, respectively, and let $r,t$ be the time and rate of Mike. Using $d = rt$, we have that $D_M = rt$ and $D_J = \left(\frac{4}{5}r\right)\left(2t\right) = \frac 85rt$. Then $13 = D_M + D_J = rt + \frac{8}{5}rt = \frac{13}{5}rt$ $\Longrightarrow rt = D_M = 5\ \mathrm{(B)}$. 

 


	\item (\textit{AMC 12 12A 2005}) Square $EFGH$ is inside the square $ABCD$ so that each side of $EFGH$ can be extended to pass through a vertex of $ABCD$. Square $ABCD$ has side length $\sqrt{50}$, $E$ is between $B$ and $H$, and $BE = 1$. What is the area of the inner square $EFGH$?

                    \begin{figure}[H]
                        \centering
                        \includegraphics[width=0.5\textwidth]{Problem/AMC_12_12A_2005_7_Problem_Figure_1.png}
                    \end{figure}
                    
$\textbf{(A) } 25 \qquad \textbf{(B) } 32 \qquad \textbf{(C) } 36 \qquad \textbf{(D) } 40 \qquad \textbf{(E) } 42$

 \textbf{Answer: }C



\textbf{Solution}


                    \begin{figure}[H]
                        \centering
                        \includegraphics[width=0.5\textwidth]{Solution/AMC_12_12A_2005_7_Solution_Figure_1.png}
                    \end{figure}
                    

 Arguable the hardest part of this question is to visualize the diagram. Since each side of $EFGH$ can be extended to pass through a vertex of $ABCD$, we realize that $EFGH$ must be tilted in such a fashion. Let a side of $EFGH$ be $x$. 

 
                    \begin{figure}[H]
                        \centering
                        \includegraphics[width=0.5\textwidth]{Solution/AMC_12_12A_2005_7_Solution_Figure_2.png}
                    \end{figure}
                    

 Notice the \href{https://artofproblemsolving.com/wiki/index.php?title=Right\_triangle}{right triangle} (in blue) with legs $1, x+1$ and \href{https://artofproblemsolving.com/wiki/index.php?title=Hypotenuse}{hypotenuse} $\sqrt{50}$. By the \href{https://artofproblemsolving.com/wiki/index.php?title=Pythagorean\_Theorem}{Pythagorean Theorem}, we have $1^2 + (x+1)^2 = (\sqrt{50})^2 \Longrightarrow (x+1)^2 = 49 \Longrightarrow x = 6$. Thus, $[EFGH] = x^2 = 36\ \mathrm{(C)}$



	\item (\textit{AMC 12 12A 2005}) Let $A$, $M$, and $C$ be digits with


 
 \[(100A+10M+C)(A+M+C) = 2005.\]


 What is $A$?


 $\textbf{(A) } 1 \qquad \textbf{(B) } 2 \qquad \textbf{(C) } 3 \qquad \textbf{(D) } 4 \qquad \textbf{(E) } 5$

 \textbf{Answer: }D



\textbf{Solution}

Clearly the two quantities are both integers, so we check the \href{https://artofproblemsolving.com/wiki/index.php?title=Prime\_factorization}{prime factorization} of $2005 = 5 \cdot 401$. It is easy to see now that $(A,M,C) = (4,0,1)$ works, so the answer is $\mathrm{(D)}$.

 


	\item (\textit{AMC 12 12A 2005}) There are two values of $a$ for which the equation $4x^2 + ax + 8x + 9 = 0$ has only one solution for $x$. What is the sum of those values of $a$?


 $\textbf{(A) } -16 \qquad \textbf{(B) } -8 \qquad \textbf{(C) } 0 \qquad \textbf{(D) } 8 \qquad \textbf{(E) } 20$

 \textbf{Answer: }A


	\item (\textit{AMC 12 12A 2005}) A wooden cube $n$ units on a side is painted red on all six faces and then cut into $n^3$ unit cubes. Exactly one-fourth of the total number of faces of the unit cubes are red. What is $n$? 


 $\textbf{(A) } 3 \qquad \textbf{(B) } 4 \qquad \textbf{(C) } 5 \qquad \textbf{(D) } 6 \qquad \textbf{(E) } 7$

 \textbf{Answer: }B



\textbf{Solution}

There are $6n^3$ sides total on the unit cubes, and $6n^2$ are painted red.

 $\dfrac{6n^2}{6n^3}=\dfrac{1}{4} \Rightarrow n=4 \rightarrow \mathrm {B}$



	\item (\textit{AMC 12 12A 2005}) How many three-digit numbers satisfy the property that the middle digit is the average of the first and the last digits?


 $\textbf{(A) } 41 \qquad \textbf{(B) } 42 \qquad \textbf{(C) } 43 \qquad \textbf{(D) } 44 \qquad \textbf{(E) } 45$

 \textbf{Answer: }E



\textbf{Solution 1}

Let the digits be $A, B, C$ so that $B = \frac {A + C}{2}$. In order for this to be an integer, $A$ and $C$ have to have the same \href{https://artofproblemsolving.com/wiki/index.php?title=Parity}{parity}. There are $9$ possibilities for $A$, and $5$ for $C$. $B$ depends on the value of both $A$ and $C$ and is unique for each $(A,C)$. Thus our answer is $9 \cdot 5 \cdot 1 = 45 \implies E$.

 



\textbf{Solution 2}

Thus, the three digits form an \href{https://artofproblemsolving.com/wiki/index.php?title=Arithmetic\_sequence}{arithmetic sequence}. 

 
            \begin{itemize}
        \item If the numbers are all the same, then there are $9$ possible three-digit numbers.
        \item If the numbers are different, then we count the number of strictly increasing arithmetic sequences between $0$ and $10$ and multiply by 2 for the decreasing ones:
            \end{itemize}
            


	\item (\textit{AMC 12 12A 2005}) A line passes through $A(1,1)$ and $B(100,1000)$. How many other points with integer coordinates are on the line and strictly between $A$ and $B$?


 $\textbf{(A) } 0 \qquad \textbf{(B) } 2 \qquad \textbf{(C) } 3 \qquad \textbf{(D) } 8 \qquad \textbf{(E) } 9$

 \textbf{Answer: }D



\textbf{Solution}

For convenience’s sake, we can transform $A$ to the origin and $B$ to $(99,999)$ (this does not change the problem). The line $AB$ has the \href{https://artofproblemsolving.com/wiki/index.php?title=Equation}{equation} $y = \frac{999}{99}x = \frac{111}{11}x$. The coordinates are integers if $11|x$, so the values of $x$ are $11, 22 \ldots 88$, with a total of $8\implies \boxed{\mathrm{(D)}}$ coordinates.

 



\textbf{Solution 2}

The slope of the line is$\frac{1000-1}{100-1}=\frac{111}{11},$so all points on the line have the form $(1+11t, 1+111t)$ for some value of $t$ (the rise is 111 and the run is 11). Such a point has integer coordinates if and only if $t$ is an integer, and the point is strictly between $A$ and $B$ if and only if $0<t<9$. Thus, there are $\boxed{8}$ points with the required property.-Paixiao

 



\textbf{Solution 3 (modular arithmetic)}

We can re-write the equation in slope-intercept form (where y is on the left side). We know that the slope is $\frac{1000 - 1}{100 - 1} = \frac{999}{99} = \frac{111}{11}$. Then, we have $y = \frac{111}{11}x - \frac{100}{11}$ which reduces to $y = \frac{111x - 100}{11}$. Now, it remains to look for values of $x$ such that $111x \equiv 1 (\mod 11)$. Since $111 \equiv 1 (\mod 11)$, there are $\boxed{8 \textbf{(D)}}$ coordinates for which this is true. ~elpianista227

 


	\item (\textit{AMC 12 12A 2005}) In the five-sided star shown, the letters $A$, $B$, $C$, $D$ and $E$ are replaced by the numbers $3$, $5$, $6$, $7$ and $9$, although not necessarily in that order. The sums of the numbers at the ends of the line segments $\overline{AB}$, $\overline{BC}$, $\overline{CD}$, $\overline{DE}$ and $\overline{EA}$ form an arithmetic sequence, although not necessarily in that order. What is the middle term of the arithmetic sequence?


 
                    \begin{figure}[H]
                        \centering
                        \includegraphics[width=0.5\textwidth]{Problem/AMC_12_12A_2005_13_Problem_Figure_1.png}
                    \end{figure}
                    
$\textbf{(A) } 9 \qquad \textbf{(B) } 10 \qquad \textbf{(C) } 11 \qquad \textbf{(D) } 12 \qquad \textbf{(E) } 13$

 \textbf{Answer: }D



\textbf{Solution 1}

$(A+B) + (B+C) + (C+D) + (D+E) + (E+A) = 2(A+B+C+D+E)$ (i.e., each number is counted twice). The sum $A + B + C + D + E$ will always be $3 + 5 + 6 + 7 + 9 = 30$, so the arithmetic sequence has a sum of $2 \cdot 30 = 60$. The middle term must be the average of the five numbers, which is $\frac{60}{5} = 12 \Longrightarrow \mathrm{(D)}$.

 



\textbf{Solution 2}

Let the terms in the arithmetic sequence be $a$, $a + d$, $a + 2d$, $a + 3d$, and $a + 4d$. We seek the middle term $a + 2d$.

 These five terms are $A + B$, $B + C$, $C + D$, $D + E$, and $E + A$, in some order. The numbers $A$, $B$, $C$, $D$, and $E$ are equal to 3, 5, 6, 7, and 9, in some order, so
 \[A + B + C + D + E = 3 + 5 + 6 + 7 + 9 = 30.\]Hence, the sum of the five terms is
 \[(A + B) + (B + C) + (C + D) + (D + E) + (E + A) = 2A + 2B + 2C + 2D + 2E = 60.\]But adding all five numbers, we also get $a + (a + d) + (a + 2d) + (a + 3d) + (a + 4d) = 5a + 10d$, so
 \[5a + 10d = 60.\]Dividing both sides by 5, we get $a + 2d = \boxed{12}$, which is the middle term. The answer is (D).

 



\textbf{Solution 3}

Not too bad with some logic and the awesome guess and check. Let $A=6$. Then let $B=7,E=5$ and $C=3,D=9$. Our arithmetic sequence is $10,11,12,13,14$ so our answer is $12 \Longrightarrow \mathrm{(D)}$.

 Solution by franzliszt

 


	\item (\textit{AMC 12 12A 2005}) On a standard die one of the dots is removed at random with each dot equally likely to be chosen. The die is then rolled. What is the probability that the top face has an odd number of dots?


 $\textbf{(A) } \frac{5}{11} \qquad \textbf{(B) } \frac{10}{21} \qquad \textbf{(C) } \frac{1}{2} \qquad \textbf{(D) } \frac{11}{21} \qquad \textbf{(E) } \frac{6}{11}$

 \textbf{Answer: }D



\textbf{Solution 1}

There are $1 + 2 + 3 + 4 + 5 + 6 = 21$ dots total. \href{https://artofproblemsolving.com/wiki/index.php?title=Casework}{Casework}:

 
            \begin{itemize}
        \item The dot is removed from an even face. There is a $\frac{2+4+6}{21} = \frac{4}{7}$ chance of this happening. Then there are 4 odd faces, giving us a probability of $\frac 47 \cdot \frac 46 = \frac{8}{21}$.
        \item The dot is removed from an odd face. There is a $\frac{1+3+5}{21} = \frac{3}{7}$ chance of this happening. Then there are 2 odd faces, giving us a probability of $\frac 37 \cdot \frac 26 = \frac{1}{7}$.
            \end{itemize}
            Thus the answer is $\frac 8{21} + \frac 17 = \frac{11}{21} \Longrightarrow \mathrm{(D)}$.

 



\textbf{Solution 2 (Alcumus)}

The probability that the top face is odd is $1/3$ if a dot is removed from an odd face, and the probability that the top face is odd is $2/3$ if a dot is removed from an even face. Because each dot has the probability $1/21$ of being removed, the top face is odd with probability$\left(\frac{1}{3}\right)\left(\frac{1+3+5}{21}\right) +\left(\frac{2}{3}\right)\left(\frac{2+4+6}{21}\right) = \frac{33}{63} = \boxed{\frac{11}{21}}$.

 


	\item (\textit{AMC 12 12A 2005}) Let $\overline{AB}$ be a diameter of a circle and $C$ be a point on $\overline{AB}$ with $2 \cdot AC = BC$. Let $D$ and $E$ be points on the circle such that $\overline{DC} \perp \overline{AB}$ and $\overline{DE}$ is a second diameter. What is the ratio of the area of $\triangle DCE$ to the area of $\triangle ABD$?


 
                    \begin{figure}[H]
                        \centering
                        \includegraphics[width=0.5\textwidth]{Problem/AMC_12_12A_2005_15_Problem_Figure_1.png}
                    \end{figure}
                    
$\textbf{(A) } \frac{1}{6} \qquad \textbf{(B) } \frac{1}{4} \qquad \textbf{(C) } \frac{1}{3} \qquad \textbf{(D) } \frac{1}{2} \qquad \textbf{(E) } \frac{2}{3}$

 \textbf{Answer: }C



\textbf{Solution 1}

Notice that the bases of both triangles are diameters of the circle. Hence the ratio of the areas is just the ratio of the heights of the triangles, or $\frac{CF}{CD}$ ($F$ is the foot of the \href{https://artofproblemsolving.com/wiki/index.php?title=Perpendicular}{perpendicular} from $C$ to $DE$).

 Call the radius $r$. Then $AC = \frac 13(2r) = \frac 23r$, $CO = \frac 13r$. Using the \href{https://artofproblemsolving.com/wiki/index.php?title=Pythagorean\_Theorem}{Pythagorean Theorem} in $\triangle OCD$, we get $\frac{1}{9}r^2 + CD^2 = r^2 \Longrightarrow CD = \frac{2\sqrt{2}}3r$. 

 Now we have to find $CF$. Notice $\triangle OCD \sim \triangle OFC$, so we can write the \href{https://artofproblemsolving.com/wiki/index.php?title=Proportion}{proportion}:

 $\frac{OF}{OC} = \frac{OC}{OD}$
$\frac{OF}{\frac{1}{3}r} = \frac{\frac{1}{3}r}{r}$
$OF = \frac 19r$By the \href{https://artofproblemsolving.com/wiki/index.php?title=Pythagorean\_Theorem}{Pythagorean Theorem} in $\triangle OFC$, we have $\left(\frac{1}{9}r\right)^2 + CF^2 = \left(\frac{1}{3}r\right)^2 \Longrightarrow CF = \sqrt{\frac{8}{81}r^2} = \frac{2\sqrt{2}}{9}r$.

 Our answer is $\frac{CF}{CD} = \frac{\frac{2\sqrt{2}}{9}r}{\frac{2\sqrt{2}}{3}r} = \frac 13 \Longrightarrow \mathrm{(C)}$.

 



\textbf{Solution 2}

Let the center of the circle be $O$.

 Note that $2 \cdot AC = BC \Rightarrow 3 \cdot AC = AB$.

 $O$ is midpoint of $AB \Rightarrow \frac{3}{2}AC = AO \Rightarrow CO = \frac{1}{3}AO \Rightarrow CO = \frac{1}{6} AB$.

 $O$ is midpoint of $DE \Rightarrow$ Area of $\triangle DCE = 2 \cdot$ Area of $\triangle DCO = 2 \cdot (\frac{1}{6} \cdot$ Area of $\triangle ABD) = \frac{1}{3} \cdot$ Area of $\triangle ABD \Longrightarrow \mathrm{(C)}$.

 



\textbf{Solution 3}

Let $r$ be the radius of the circle. Note that $AC+BC = 2r$ so $AC = \frac{2}{3}r$.

 By \href{https://artofproblemsolving.com/wiki/index.php?title=Power\_of\_a\_Point\_Theorem}{Power of a Point Theorem}, $CD^2= AC \cdot BC = 2\cdot AC^2$, and thus $CD = \sqrt{2} \cdot AC = \frac{2\sqrt{2}}{3}r$

 Then the area of $\triangle ABD$ is $\frac{1}{2} AB \cdot CD = \frac{2\sqrt{2}}{3}r^2$. Similarly, the area of $\triangle DCE$ is $\frac{1}{2}(r-AC) \cdot 2 \cdot CD = \frac{2\sqrt{2}}{9}r^2$, so the desired ratio is $\frac{\frac{2\sqrt{2}}{9}r^2}{\frac{2\sqrt{2}}{3}r^2} = \frac{1}{3} \Longrightarrow \mathrm{(C)}$




\textbf{Solution 4}


                    \begin{figure}[H]
                        \centering
                        \includegraphics[width=0.5\textwidth]{Solution/AMC_12_12A_2005_15_Solution_Figure_1.png}
                    \end{figure}
                    

 Let the center of the circle be $O$.Without loss of generality, let the radius of the circle be equal to $3$. Thus, $AO=3$ and $OB=3$. As a consequence of $2(AC)=BC$, $AC=2$ and $CO=1$. Also, we know that $DO$ and $OE$ are both equal to $3$ due to the fact that they are both radii. Thus from the Pythagorean Theorem, we have DC being equal to $\sqrt{3^2-1^2}$ or $2\sqrt{2}$. Now we know that the area of $[ABD]$ is equal to $\frac{(3+2+1)(2\sqrt{2})}{2}$ or $6\sqrt{2}$. Know we need to find the area of $[DCE]$. By simple inspection $[COD]$ $\cong$ $[HOE]$ due to angles being equal and CPCTC. Thus $HE=2\sqrt{2}$ and $OH=1$. Know we know the area of $[CHE]=\frac{(1+1)(2\sqrt{2})}{2}$ or $2\sqrt{2}$. We also know that the area of $[OHE]=\frac{(1)(2\sqrt{2})}{2}$ or $\sqrt{2}$. Thus the area of $[COE]=2\sqrt{2}-\sqrt{2}$ or $\sqrt{2}$. We also can calculate the area of $[DOC]$ to be $\frac{(1)(2\sqrt{2})}{2}$ or $\sqrt{2}$. Thus $[DCE]$ is equal to $[COE]$ + $[DOC]$ or $\sqrt{2}+\sqrt{2}$ or $2\sqrt{2}$. The ratio between $[DCE]$ and  $[ABD]$ is equal to $\frac{2\sqrt{2}}{6\sqrt{2}}$ or $\frac{1}{3}$ $\Longrightarrow \mathrm{(C)}$.

 



\textbf{Solution 5}

We will use the shoelace formula. Our origin is the center of the circle. Denote the ordered pair for $D (x,y)$, and notice how $E$ is a 180 degree rotation of $D$, using the rotation matrix formula we get $E = (-x,-y)$. WLOG say that this circle has radius $3$. We can now find points $A$, $B$, and $C$ which are $(-3,0)$, $(3,0)$, and $(-1,0)$ respectively.By shoelace the area of $CED$ is $Y$, and the area of $ADB$ is $3Y$. Using division we get that the answer is $\mathrm{(C)}$.

 



\textbf{Solution 6 (Mass Points)}


                    \begin{figure}[H]
                        \centering
                        \includegraphics[width=0.5\textwidth]{Solution/AMC_12_12A_2005_15_Solution_Figure_2.png}
                    \end{figure}
                    

 We set point $A$ as a mass of 2. This means that point $B$ has a mass of $1$ since $2\times{AC} = 1\times{BC}$. This implies that point $C$ has a mass of $2+1 = 3$ and the center of the circle has a mass of $3+1 = 4$. After this, we notice that points $D$ and $E$ both must have a mass of $2$ since $2+2 = 4$ and they are both radii of the circle.

 To find the ratio of the areas, we do the reciprocal of the multiplication of the mass points of DCE and the multiplication of ABD divided by each other. Which is simply $\frac{3\times2\times2}{2\times2\times1}$ which is $\boxed{\frac{1}{3}}$ (the reciprocal of 3)

 -Brudder

 



\textbf{Solution 7 (Slight Trigonometry)}


                    \begin{figure}[H]
                        \centering
                        \includegraphics[width=0.5\textwidth]{Solution/AMC_12_12A_2005_15_Solution_Figure_3.png}
                    \end{figure}
                    

 Let the center of the circle be $O$.The area of $ADB$ is $\frac{1}{2} \cdot 6x \cdot 2x\sqrt{2}$ = $6x^2\sqrt{2}$. The area of $DCE$ is $\frac{1}{2} \cdot DC \cdot DE \cdot$ sin $CDE$. We find sin $CDE$ is $\frac{OC}{OD}$ = $\frac{1}{3}$. Substituting $DC = 2x\sqrt{2}$ and $DE = 6x$, we get $[DCE]$ = $\frac{1}{2} \cdot 2x\sqrt{2} \cdot 6x \cdot \frac{1}{3}$ = $2x^2\sqrt{2}$. Hence, the ratio between the areas of $DCE$ and $ADB$ is equal to $\frac{2x^2\sqrt{2}}{6x^2\sqrt{2}}$ or $\frac{1}{3}$ = $\boxed{(C)}$.

 



\textbf{Solution 8}

In my opinion, the solution below is the easiest and quickest. 

 Since both $DE$ and $AB$ are diameters, they intersect at the center of the circle. Call this center $O$. WLOG, let $AC=2, CO=1,OB=3$. Call the point where the extension of $DC$ hits the circle $F$. Notice that $\angle DCB = 90^\circ$. This implies that $DC=CF$. WOLG, let $DC=CF=1$. Then, $[DCE]=\frac 12 bh = \frac 12 * DF*CO = \frac 12 *2*1=1$ and $[ABD]=\frac 12 bh = \frac 12 *AB*CD = \frac 12 *6*1 = 3$. Thus, the answer is $\frac{1}{3}$ = $\boxed{(C)}$.

 Solution by franzliszt

 



\textbf{Solution 9 (Slick construction)}


                    \begin{figure}[H]
                        \centering
                        \includegraphics[width=0.5\textwidth]{Solution/AMC_12_12A_2005_15_Solution_Figure_4.png}
                    \end{figure}
                    Let $X$ be the reflection of $C$ over the center $O.$ Since $\triangle DXO\cong\triangle ECO$ by SAS, it follows that the area of $\triangle DEC$ is equal to the area of $\triangle DCX.$ However, we know that 
 \[BC = 2AC\implies CO=\frac{1}{2}AC\implies CX=AC,\]so the ratio of the area of $\triangle DEC$ to the area of $\triangle ABD$ is $\frac{\frac{AC\cdot DC}{2}}{\frac{3AC\cdot DC}{2}}=\frac{1}{3}\implies \boxed{(C)}$



	\item (\textit{AMC 12 12A 2005}) Three circles of radius $s$ are drawn in the first quadrant of the $xy$-plane. The first circle is tangent to both axes, the second is tangent to the first circle and the $x$-axis, and the third is tangent to the first circle and the $y$-axis. A circle of radius $r > s$ is tangent to both axes and to the second and third circles. What is $r/s$?


 
                    \begin{figure}[H]
                        \centering
                        \includegraphics[width=0.5\textwidth]{Problem/AMC_12_12A_2005_16_Problem_Figure_1.png}
                    \end{figure}
                    
$\textbf{(A) } 5 \qquad \textbf{(B) } 6 \qquad \textbf{(C) } 8 \qquad \textbf{(D) } 9 \qquad \textbf{(E) } 10$

 \textbf{Answer: }D



\textbf{Solution}


                    \begin{figure}[H]
                        \centering
                        \includegraphics[width=0.5\textwidth]{Solution/AMC_12_12A_2005_16_Solution_Figure_1.png}
                    \end{figure}
                    

 Set $s =1$ so that we only have to find $r$. Draw the segment between the center of the third circle and the large circle; this has length $r+1$. We then draw the \href{https://artofproblemsolving.com/wiki/index.php?title=Radius}{radius} of the large circle that is perpendicular to the \href{https://artofproblemsolving.com/wiki/index.php?title=X-axis}{x-axis}, and draw the perpendicular from this radius to the center of the third circle. This gives us a \href{https://artofproblemsolving.com/wiki/index.php?title=Right\_triangle}{right triangle} with legs $r-3,r-1$ and \href{https://artofproblemsolving.com/wiki/index.php?title=Hypotenuse}{hypotenuse} $r+1$. The \href{https://artofproblemsolving.com/wiki/index.php?title=Pythagorean\_Theorem}{Pythagorean Theorem} yields:

 $(r-3)^2 + (r-1)^2 = (r+1)^2$
$r^2 - 10r + 9 = 0$
$r = 1, 9$Quite obviously $r > 1$, so $r = 9$.

 


	\item (\textit{AMC 12 12A 2005}) A unit cube is cut twice to form three triangular prisms, two of which are congruent, as shown in Figure $1$. The cube is then cut in the same manner along the dashed lines shown in Figure $2$. This creates nine pieces. What is the volume of the piece that contains vertex $W$?


 
                    \begin{figure}[H]
                        \centering
                        \includegraphics[width=0.5\textwidth]{Problem/AMC_12_12A_2005_17_Problem_Figure_1.png}
                    \end{figure}
                    
$\textbf{(A) } \frac{1}{12} \qquad \textbf{(B) } \frac{1}{9} \qquad \textbf{(C) } \frac{1}{8} \qquad \textbf{(D) } \frac{1}{6} \qquad \textbf{(E) } \frac{1}{4}$

 \textbf{Answer: }A



\textbf{Solution}

It is a \href{https://artofproblemsolving.com/wiki/index.php?title=Pyramid}{pyramid} with height $1$ and base area $\frac{1}{4}$, so using the formula for the volume of a pyramid, $\frac{1}{3} \cdot \left(\frac{1}{4}\right) \cdot (1) = \frac {1}{12} \Rightarrow \boxed{(\mathrm {A})}$.

 


	\item (\textit{AMC 12 12A 2005}) Call a number "prime-looking" if it is composite but not divisible by $2$, $3$, or $5$. The three smallest prime-looking numbers are $49$, $77$, and $91$. There are $168$ prime numbers less than $1000$. How many prime-looking numbers are there less than $1000$?


 $\textbf{(A) } 100 \qquad \textbf{(B) } 102 \qquad \textbf{(C) } 104 \qquad \textbf{(D) } 106 \qquad \textbf{(E) } 108$

 \textbf{Answer: }A



\textbf{Solution}

The given states that there are $168$ prime numbers less than $1000$, which is a fact we must somehow utilize. Since there seems to be no easy way to directly calculate the number of "prime-looking" numbers, we can apply \href{https://artofproblemsolving.com/wiki/index.php?title=Complementary\_counting}{complementary counting}. We can split the numbers from $1$ to $1000$ into several groups: $\{1\},$ $\{\mathrm{numbers\ divisible\ by\ 2 = S_2}\},$ $\{\mathrm{numbers\ divisible\ by\ 3 = S_3}\},$ $\{\mathrm{numbers\ divisible\ by\ 5 = S_5}\}, \{\mathrm{primes\ not\ including\ 2,3,5}\},$ $\{\mathrm{prime-looking}\}$. Hence, the number of prime-looking numbers is $1000 - (168-3) - 1 - |S_2 \cup S_3 \cup S_5|$ (note that $2,3,5$ are primes).

 We can calculate $S_2 \cup S_3 \cup S_5$ using the \href{https://artofproblemsolving.com/wiki/index.php?title=Principle\_of\_Inclusion-Exclusion}{Principle of Inclusion-Exclusion}: (the values of $|S_2| \ldots$ and their intersections can be found quite easily)

 $|S_2 \cup S_3 \cup S_5| = |S_2| + |S_3| + |S_5| - |S_2 \cap S_3| - |S_3 \cap S_5| - |S_2 \cap S_5| + |S_2 \cap S_3 \cap S_5|$
$= 500 + 333 + 200 - 166 - 66 - 100 + 33 = 734$Substituting, we find that our answer is $1000 - 165 - 1 - 734 = 100 \Longrightarrow \mathrm{(A)}$.

 


	\item (\textit{AMC 12 12A 2005}) A faulty car odometer proceeds from digit $3$ to digit $5$, always skipping the digit $4$, regardless of position. For example, after traveling one mile the odometer changed from $000039$ to $000050$. If the odometer now reads $002005$, how many miles has the car actually traveled?


 $\textbf{(A) } 1404 \qquad \textbf{(B) } 1462 \qquad \textbf{(C) } 1604 \qquad \textbf{(D) } 1605 \qquad \textbf{(E) } 1804$

 \textbf{Answer: }B



\textbf{Solution 1}

We find the number of numbers with a $4$ and subtract from $2005$. Quick counting tells us that there are $200$ numbers with a 4 in the hundreds place, $200$ numbers with a 4 in the tens place, and $201$ numbers with a 4 in the units place (counting $2004$). Now we apply the \href{https://artofproblemsolving.com/wiki/index.php?title=Principle\_of\_Inclusion-Exclusion}{Principle of Inclusion-Exclusion}. There are $20$ numbers with a 4 in the hundreds and in the tens, and $20$ for both the other two \href{https://artofproblemsolving.com/wiki/index.php?title=Intersection}{intersections}. The intersection of all three sets is just $2$. So we get:

 $2005-(200+200+201-20-20-20+2) = 1462 \Longrightarrow \mathrm{(B)}$



\textbf{Solution 2}

Alternatively, consider that counting without the number $4$ is almost equivalent to counting in base $9$; only, in base $9$, the number $9$ is not counted. Since $4$ is skipped, the symbol $5$ represents $4$ miles of travel, and we have traveled $2004_9$ miles. By    base conversion, $2004_9=9^3(2)+9^0(4)=729(2)+1(4)=1458+4=\boxed{1462}$




\textbf{Solution 3}

Since any numbers containing one or more $4$s were skipped, we need only to find the numbers that don't contain a $4$ at all. First we consider $1$ - $1999$. Single digits are simply four digit numbers with a zero in all but the ones place (this concept applies to double and triple digits numbers as well). From $1$ - $1999$, we have $2$ possibilities for the thousands place, and $9$ possibilities for the hundreds, tens, and ones places. This is $2 \cdot 9 \cdot 9 \cdot 9-1$ possibilities (because $0000$ doesn't count) or $1457$ numbers. From $2000$ - $2005$ there are $6$ numbers, $5$ of which don't contain a $4$. Therefore the total is $1457 + 5$, or $1462$ $\Rightarrow$ $\boxed{\text{B}}$.

 



\textbf{Solution 4}

We seek to find    the amount of numbers that contain at least one $4,$ and subtract this number from $2005.$ 

 We can simply apply casework to this problem.

 The amount of numbers with at least one $4$ that are one or two digit numbers are $4,14,24,34,40-49,54,\cdots,94$ which gives $19$ numbers.

 The amount of three digit numbers with at least one $4$ is $8*19+100=252.$ 

 The amount of four digit numbers with at least one $4$ is $252+1+19=272$

 This, our answer is $2005-19-252-272=1462,$ or $\boxed{B}.$




\textbf{Solution 5 (Super fast)}

This is very analogous to base $9$. But, in base $9$, we don't have a $9$. So, this means that these are equal except for that base 9 will be one more than the operation here. $2005_9 = 5+0+0+1458 = 1463$. $1463 - 1 = 1462$

 Therefore, our answer is $\boxed {1462}$



	\item (\textit{AMC 12 12A 2005}) For each $x$ in $[0,1]$, define


 
 \[f(x) = \begin{cases}  2x, \qquad\qquad \mathrm{if} \quad 0 \leq x \leq \frac{1}{2}\\  2-2x, \qquad \mathrm{if} \quad \frac{1}{2} < x \leq 1.  \end{cases}\]


 Let $f^{[2]}(x) = f(f(x))$, and $f^{[n + 1]}(x) = f^{[n]}(f(x))$ for each integer $n \geq 2$. For how many values of $x$ in $[0,1]$ is $f^{[2005]}(x) = 1/2$?


 $\textbf{(A) } 0 \qquad \textbf{(B) } 2005 \qquad \textbf{(C) } 4010 \qquad \textbf{(D) } 2005^2 \qquad \textbf{(E) } 2^{2005}$

 \textbf{Answer: }E



\textbf{Solution 1}

For the two functions $f(x)=2x,0\le x\le \frac{1}{2}$ and $f(x)=2-2x,\frac{1}{2}< x\le 1$,as long as $f(x)$ is between $0$ and $1$, $x$ will be in the right domain, so we don't need to worry about the domain of $x$.  

 Also, every time we change $f(x)$, the expression for the final answer in terms of $x$ will be in a different form (although they'll all satisfy the final equation), so we get a different starting value of $x$.  Every time we have two choices for $f(x$) and altogether we have to choose $2005$ times. Thus, $2^{2005}\Rightarrow\boxed{E}$.

 Note: the values of x that satisfy $f^{[n]}(x) = \frac {1}{2}$ are $\frac{1}{2^{n+1}}$, $\frac{3}{2^{n+1}}$, $\frac{5}{2^{n+1}}$, $\cdots$ ,$\frac{2^{n+1}-1}{2^{n+1}}$.

 



\textbf{Solution 2 (Engineers Induction)}

We are given that $f^{[2005]}(x) = \frac {1}{2}$. Thus,  $f(f^{[2004]}(x))=\frac{1}{2}$. Let $f^{[2004]}(x)$ be equal to $y$. Thus $f(y)=\frac{1}{2}$ or $y=\frac{1}{4}$ or $\frac{3}{4}$. Now we know $f^{[2004]}(x)$ is equal to $\frac{1}{4}$ or $\frac{3}{4}$. Now we know that $f(f^{[2003]}(x))=\frac{1}{4}$ or $\frac{3}{4}$. Now we solve for $f^{[2003]}(x)$ and let  $f^{[2003]}(x)=z$. Thus $f(z)$ is equal to $\frac{1}{8}$,$\frac{7}{8}$,$\frac{5}{8}$,and $\frac{3}{8}$. As we see, $f^{[2005]}(x)$ has 1 solution, $f^{[2004]}(x)$ has 2 solutions, and $f^{[2003]}(x)$ has 4 solutions. Thus for each iteration we double the number of possible solutions. There are $2005$ iterations and thus the number of solutions is $2^{2005}$ $\Rightarrow\boxed{E}$



	\item (\textit{AMC 12 12A 2005}) How many ordered triples of integers $(a,b,c)$, with $a \geq 2$, $b \geq 1$, and $c \geq 0$, satisfy both $\log_{a}b = c^{2005}$ and $a + b + c = 2005$?


 $\textbf{(A) } 0 \qquad \textbf{(B) } 1 \qquad \textbf{(C) } 2 \qquad \textbf{(D) } 3 \qquad \textbf{(E) } 4$

 \textbf{Answer: }C



\textbf{Solution}

$a^{c^{2005}} = b$

 \href{https://artofproblemsolving.com/wiki/index.php?title=Casework}{Casework} upon $c$:

 
            \begin{itemize}
        \item $c = 0$: Then $a^0 = b \Longrightarrow b = 1$. Thus we get $(2004,1,0)$.
        \item $c = 1$: Then $a^1 = b \Longrightarrow a = b$. Thus we get $(1002,1002,1)$.
        \item $c \ge 2$: Then the exponent of $a$ becomes huge, and since $a \ge 2$ there is no way we can satisfy the second condition. Hence we have two ordered triples $\mathrm{(C)}$.
            \end{itemize}
            


	\item (\textit{AMC 12 12A 2005}) A rectangular box $P$ is inscribed in a sphere of radius $r$. The surface area of $P$ is $384$, and the sum of the lengths of its $12$ edges is $112$. What is $r$?


 $\textbf{(A) } 8 \qquad \textbf{(B) } 10 \qquad \textbf{(C) } 12 \qquad \textbf{(D) } 14 \qquad \textbf{(E) } 16$

 \textbf{Answer: }B



\textbf{Solution 1}

Box P has dimensions $l$, $w$, and $h$. Its surface area is 
 \[2lw+2lh+2wh=384,\] and the sum of all its edges is 
 \[l + w + h = \dfrac{4l+4w+4h}{4} = \dfrac{112}{4} = 28.\]

 The diameter of the sphere is the space diagonal of the prism, which is 
 \[\sqrt{l^2 + w^2 +h^2}.\]Notice that 
 \begin{align*}(l + w + h)^2 - (2lw + 2lh + 2wh) &= l^2 + w^2 + h^2 \\ &= 784 - 384 \\ &= 400,\end{align*}
 so the diameter is
 \[\sqrt{l^2 + w^2 +h^2} = \sqrt{400} = 20.\] The radius is half of the diameter, so
 \[r=\frac{20}{2} = \boxed{\textbf{(B) } 10}.\]




\textbf{Solution 2 (non-rigorous)}

As in the previous solution, we have that $2lw+2lh+2wh=384$ and $l + w + h = \dfrac{4l+4w+4h}{4} = \dfrac{112}{4} = 28$, and the diameter of the sphere is the space diagonal of the prism, $\sqrt{l^2 + w^2 + h^2}$.

 Now, since this is competition math, we only need to find the space diagonal of any one box that fits the requirements, so assume that $h=0$. (This essentially means that we have an infinitesimally thin box.) We now have that $2lw = 384$ and $l + w = 28$, and we are solving for $\sqrt{l^2 + w^2}$. Because 
 \begin{align*}(l + w)^2 - 2lw &= l^2 + 2lw + w^2 - 2lw \\ &= l^2 + w^2,\end{align*}
 this means that 
 \[l^2 + w^2 = 28^2 - 384 = 400,\] so the space diagonal is $\sqrt{400} = 20$. Since the diameter of the sphere is $20$, the radius is $\boxed{\textbf{(B) } 10}$. ~\href{https://artofproblemsolving.com/wiki/index.php?title=User:Emerald\_block}{emerald\_block}

 


	\item (\textit{AMC 12 12A 2005}) Two distinct numbers $a$ and $b$ are chosen randomly from the set $\{2, 2^2, 2^3, \ldots, 2^{25}\}$. What is the probability that $\log_{a}b$ is an integer?


 $\textbf{(A) } \frac{2}{25} \qquad \textbf{(B) } \frac{31}{300} \qquad \textbf{(C) } \frac{13}{100} \qquad \textbf{(D) } \frac{7}{50} \qquad \textbf{(E) } \frac{1}{2}$

 \textbf{Answer: }B



\textbf{Solution}

Let $\log_{a}b = z$, so $a^z = b$. Define $a = 2^x$, $b = 2^y$; then $\left(2^x\right)^z = 2^{xz}= 2^y$, so $x \mid y$. Here we can just make a table and count the number of values of $y$ per value of $x$. The largest possible value of $x$ is $12$, and so we get $\sum_{x=1}^{12} \left\lfloor\frac{25}{x}-1\right\rfloor = 24 + 11 + 7 + 5 + 4 + 3 + 2 + 2 + 1 + 1 + 1 + 1 = 62$. 

 The total number of ways to pick two distinct numbers (where the order matters) is $\frac{25!}{(25-2)!}= 25 \cdot 24 = 600$, so we get a probability of $\frac{62}{600} = \boxed{\textbf{(B) }\frac{31}{300}}$.

 


	\item (\textit{AMC 12 12A 2005}) Let $P(x) = (x - 1)(x - 2)(x - 3)$. For how many polynomials $Q(x)$ does there exist a polynomial $R(x)$ of degree $3$ such that $P(Q(x)) = P(x) \cdot R(x)$?


 $\textbf{(A) } 19 \qquad \textbf{(B) } 22 \qquad \textbf{(C) } 24 \qquad \textbf{(D) } 27 \qquad \textbf{(E) } 32$

 \textbf{Answer: }B



\textbf{Solution 1}

We can write the problem as

 $P(Q(x))=(Q(x)-1)(Q(x)-2)(Q(x)-3)=P(x)\cdot R(x)=(x-1)(x-2)(x-3)\cdot R(x)$.

 
Since $\deg P(x) = 3$ and $\deg R(x) = 3$, $\deg P(x)\cdot R(x) = 6$. Thus, $\deg P(Q(x)) = 6$, so $\deg Q(x) = 2$.

 <div style="text-align:center;">$P(Q(1))=(Q(1)-1)(Q(1)-2)(Q(1)-3)=P(1)\cdot R(1)=0,$
$P(Q(2))=(Q(2)-1)(Q(2)-2)(Q(2)-3)=P(2)\cdot R(2)=0,$
$P(Q(3))=(Q(3)-1)(Q(3)-2)(Q(3)-3)=P(3)\cdot R(3)=0.$
</div>Hence, we conclude $Q(1)$, $Q(2)$, and $Q(3)$ must each be $1$, $2$, or $3$. Since a quadratic is uniquely determined by three points, there can be $3*3*3 = 27$ different quadratics $Q(x)$ after each of the values of $Q(1)$, $Q(2)$, and $Q(3)$ are chosen.

 
However, we have included polynomials $Q(x)$ which are linear, rather than quadratic. Namely,

 <div style="text-align:center;">$Q(1)=Q(2)=Q(3)=1 \Rightarrow Q(x)=1,$
$Q(1)=Q(2)=Q(3)=2 \Rightarrow Q(x)=2,$
$Q(1)=Q(2)=Q(3)=3 \Rightarrow Q(x)=3,$
$Q(1)=1, Q(2)=2, Q(3)=3 \Rightarrow Q(x)=x,$
$Q(1)=3, Q(2)=2, Q(3)=1 \Rightarrow Q(x)=4-x.$
</div>Clearly, we could not have included any other constant functions. For any linear function, we have $2\cdot Q(2) = Q(1) + Q(3)$ because $Q(2)$ is the $y$-coordinate of the midpoint of $(1, Q(1))$ and $(3, Q(3))$. So we have not included any other linear functions. Therefore, the desired answer is $27 - 5 = \boxed{\textbf{(B) }22}$.

 



\textbf{Solution 2}

We see that 
 \[P(Q(x))=(Q(x)-1)(Q(x)-2)(Q(x)-3)=P(x)\cdot R(x)=(x-1)(x-2)(x-3)\cdot R(x).\]Therefore, $P(x) | P(Q(x))$.  Since $\deg Q = 2,$ we must have $x-1, x-2, x-3$ divide $P(Q(x))$. So, we pair them off with one of $Q(x)-1, Q(x)-2,$ and $Q(x)-3$ to see that there are $3!+3 \cdot 2 \cdot \binom{3}{2} = 24$ without restrictions. (Note that this count was made by pairing off linear factors of $P(x)$ with $Q(x)-1, Q(x)-2,$ and $Q(x)-3$, and also note that the degree of $Q$ is $2$.) However, we have two functions which are constant, which are $Q(x) = x$ and $Q(x) = 4-x$. So we subtract $2$ to get a final answer of $\boxed{22} \implies \boxed{B}$. 

 



\textbf{Solution 3 (fastest)}

Obviously, $\deg Q = 2$ and the RHS has roots $1, 2, 3$. Thus, the only requirement is that $Q$ is a quadratic that maps each of the numbers $1, 2, 3$ to one of ${1, 2, 3}.$ It suffices to count the number of such mappings, as this uniquely determines $Q$ through polynomial interpolation. Thus there are $3^3$ possibilities. But we have three functions that are constant (i.e. that give $Q(1) =Q(2) = Q(3)$), and $2$ that are linear (i.e. that map $\{1, 2, 3\} \mapsto \{3, 2, 1\}$ or $\{1, 2, 3\} \mapsto \{1, 2, 3\}$). Hence there are $3^3-3-2 = \boxed{22}$ solutions, so $\boxed{B}$.

 \textbf{Non-Solution 4 (guessing, only if desperate!)} As $R(x)$ must have degree $3$, let its roots be $r_1$, $r_2$, and $r_3$, giving
 \begin{align*}P(Q(x)) &= (Q(x)-1)(Q(x)-2)(Q(x)-3) \\ &= P(x)\cdot R(x) \\ &= (x-1)(x-2)(x-3) \cdot R(x) \\ &= \left(x-1\right)\left(x-2\right)\left(x-3\right)\left(x-r_1\right)\left(x-r_2\right)\left(x-r_3\right).\end{align*}


 If we now guess that $Q(x)$ is a quadratic polynomial, it follows that e.g. $\left(Q(x)-1\right)$ must be equal to a product of $2$ factors from the expression $\left(x-1\right)\left(x-2\right)\left(x-3\right)\left(x-r_1\right)\left(x-r_2\right)\left(x-r_3\right)$. As this expression contains $6$ factors in total, the number of ways of choosing $2$ factors from it, where the order matters, would be $6!/4! = 30$.

 While none of the answer choices are $30$, we can now eliminate $\mathrm {(E) }$ as $32 > 30$. Looking for answers that are similar, we further observe that $22$, $27$, or $32$ would have been derived from a $27$ in the actual solution, to which $5$ was possibly added or subtracted. We can therefore guess, if sufficiently desperate, that the answer should be either $\mathrm {(B) } \ 22$ or $\mathrm {(D) } \ 27$.

 On the other hand, if you're desperate for time, it's probably a better idea to go back and review Problems $1$ to $20$, rather than guessing on Problem $24$!

 


	\item (\textit{AMC 12 12A 2005}) Let $S$ be the set of all points with coordinates $(x,y,z)$, where $x$, $y$, and $z$ are each chosen from the set $\{0, 1, 2\}$. How many equilateral triangles have all their vertices in $S$?


 $\textbf{(A) } 72 \qquad \textbf{(B) } 76 \qquad \textbf{(C) } 80 \qquad \textbf{(D) } 84 \qquad \textbf{(E) } 88$

 \textbf{Answer: }C



\textbf{Solution 1}

For this solution, we will just find as many equilateral triangles as possible, until it becomes intuitive that there are no more size of triangles left.

 First, we observe that we can form an equilateral triangle with vertices in $S$ by taking any point in $S$ and connecting it to the $2$ adjacent points. This triangle will have a side length of $\sqrt{2}$; a quick further examination of this cube will show us that this is the only possible side length (the red triangle in the diagram below). Each of these triangles is determined by one vertex of the cube, so in one cube we have $8$ equilateral triangles. We have $8$ unit cubes, as well as the entire $2 \times 2 \times 2$ cube (giving the green triangle in the diagram), for a total of $8+1 = 9$ cubes, and thus $9 \cdot 8 = 72$ equilateral triangles. 

 
                    \begin{figure}[H]
                        \centering
                        \includegraphics[width=0.5\textwidth]{Solution/AMC_12_12A_2005_25_Solution_Figure_1.png}
                    \end{figure}
                    

 

 (Note that connecting the centers of the faces will actually give an octahedron, not a cube, because it only has $6$ vertices.)

 Now we look for any further equilateral triangles. Connecting the midpoints of $3$ non-adjacent, non-parallel edges indeed gives us more equilateral triangles (e.g. the blue triangle in the diagram below).  Notice that picking these $3$ edges leaves $2$ vertices alone (labelled A and B in the diagram), and that picking any $2$ opposite vertices determines $2$ equilateral triangles. Hence there are $\frac{8 \cdot 2}{2} = 8$ of these equilateral triangles, so adding these to the triangles already found above gives a total of $72+8 = \boxed{\textbf{(C) }80}$.

 
                    \begin{figure}[H]
                        \centering
                        \includegraphics[width=0.5\textwidth]{Solution/AMC_12_12A_2005_25_Solution_Figure_2.png}
                    \end{figure}
                    




\textbf{Solution 2}

The three-dimensional distance formula shows that the side length of the equilateral triangle must be $\sqrt{d_x^2 + d_y^2 + d_z^2}$ with $\left\lvert d_x\right\rvert,\left\lvert d_y\right\rvert,\left\lvert d_z\right\rvert \in \{0-0,1-0,1-1,2-0,2-1,2-2\} = \{0,1,2\}$, so the possible side lengths are
 \begin{align*}&\sqrt{0^2+0^2+1^2}=\sqrt{1},\ \sqrt{0^2+1^2+1^2}=\sqrt{2},\ \sqrt{1^2+1^2+1^2}=\sqrt{3}, \\ &\sqrt{0^2+0^2+2^2}=\sqrt{4},\ \sqrt{0^2+1^2+2^2}=\sqrt{5},\ \sqrt{1^2+1^2+2^2}=\sqrt{6}, \\ &\sqrt{0^2+2^2+2^2}=\sqrt{8},\ \sqrt{1^2+2^2+2^2}=\sqrt{9},\ \sqrt{2^2+2^2+2^2}=\sqrt{12}.\end{align*}


 Some casework shows that $\sqrt{2}$, $\sqrt{6}$, and $\sqrt{8}$ are the only lengths that work, after which we can complete the problem using the same counting argument as in Solution 1.

 


\end{enumerate}